\chapter{SIMD Acceleration}
\label{ch:simd}

The DFS matcher's inner work is dominated by edge filtering:
checking \code{rel\_type\_tag}, \code{entity\_type\_tags[dst]},
timestamp inclusion, and predicate evaluation for millions of
candidate edges per hunt. Three of those four checks are
byte-level comparisons that vectorise cleanly. GraphHunter has
two SIMD paths:

\begin{enumerate}
    \item \textbf{libgraphmatch} --- a C++ library linked via FFI
        when the \code{simd} feature is compiled and the host CPU
        reports AVX2 support. Peak throughput around 3$\times$ the
        Rust scalar matcher.
    \item \textbf{Pure-Rust \code{simd\_rust}} (P2-D) --- a
        \code{std::arch::x86\_64} implementation that uses AVX2
        intrinsics without leaving the Rust toolchain. Currently
        implements one primitive, \textsc{Intersect-Sorted-U32},
        with more on the way.
\end{enumerate}

\section{libgraphmatch FFI bridge}

\begin{lstlisting}[style=rust]
// graph_hunter_core/src/simd_matcher.rs (abridged)
unsafe extern "C" {
    fn gm_graph_new() -> *mut c_void;
    fn gm_graph_add_node(g: *mut c_void, id: u32, et: u8);
    fn gm_graph_add_edge(g: *mut c_void, src: u32, dst: u32, rt: u8, ts: i64);
    fn gm_graph_finalize(g: *mut c_void);
    fn gm_matcher_run(m: *mut c_void, max_results: u32, min_score: f64)
         -> *mut c_void;
    fn gm_results_free(r: *mut c_void);
    fn gm_graph_free(g: *mut c_void);
}
\end{lstlisting}

The Rust side wraps the \code{*mut c\_void} handles in
\code{GmGraph} (RAII, drop calls \code{gm\_graph\_free}) and caches
the resulting graph alongside a \code{HashMap<String, u32>} that
maps string IDs to the C++ node indexes (see
\code{CachedSimdGraph} in \filepath{simd\_matcher.rs}). The cache
survives across hunts on the same session; it's invalidated
whenever the graph is mutated.

\textsc{Search-Temporal-Pattern-Simd} is therefore a short
function:

\begin{algorithm}
\caption{\textsc{Search-Temporal-Pattern-Simd}}\label{alg:search-simd}
\KwIn{graph $G$, hypothesis $H$, cap $\mathrm{max}$}
\If{runtime AVX2 not detected or $H$ has predicates}{
    \Return \textsc{Search-Temporal-Pattern-Smart}($G, H, \ldots$)
    \Comment*[r]{fallback}
}
$C \gets $ \code{G.cached\_simd\_graph()} \Comment*[r]{build once, reuse}
$r \gets $ \code{gm\_matcher\_run($C$, max, 0.0)}\;
\Return \code{translate\_node\_ids($r$, $C$.\text{string\_map})}\;
\end{algorithm}

The fallback to the Rust smart matcher is unavoidable because
libgraphmatch does not evaluate metadata predicates
(\code{\{key="val"\}}) --- those require
\code{MetadataStore} access which the C++ side cannot do
without shipping the raw bytes across FFI.

\section{Pure-Rust AVX2 intersection}

The P2-D refactor began replacing libgraphmatch one primitive at a
time. The first one shipped is
\code{simd\_rust::intersect\_sorted\_u32\_count}: given two
ascending-sorted deduplicated \code{\&[u32]} slices, return the
count of common elements. This is the primitive that
libgraphmatch uses for neighbor intersection during SIMD pattern
expansion, and it's the cheapest primitive to prove safe in
isolation.

\begin{algorithm}
\caption{\textsc{Intersect-Sorted-U32-Count}}\label{alg:intersect}
\KwIn{sorted slices $A$, $B$ of \code{u32}, both ascending and
      deduplicated}
\KwOut{count of common elements}
\If{$|A| < 32$ or $|B| < 32$}{\Return \textsc{Scalar-Merge}($A, B$)\;}
\If{AVX2 not detected at runtime}{\Return \textsc{Scalar-Merge}($A, B$)\;}
\tcp{short side is $S$, long side is $L$}
\lIf{$|A| \leq |B|$}{$(S, L) \gets (A, B)$}\lElse{$(S, L) \gets (B, A)$}
$j \gets 0$; $\mathrm{count} \gets 0$\;
\ForEach{$s \in S$}{
    \While{$j + 8 \leq |L|$}{
        $\mathrm{chunk} \gets $ AVX2 load $L[j..j+8]$\;
        $\mathrm{needle} \gets $ AVX2 broadcast $s$ to 8 lanes\;
        $\mathrm{cmp} \gets $ \code{\_mm256\_cmpeq\_epi32}(chunk, needle)\;
        $\mathrm{mask} \gets $ \code{\_mm256\_movemask\_epi8}(cmp)\;
        \If{$\mathrm{mask} \neq 0$}{
            $\mathrm{count} \gets \mathrm{count} + 1$;
            \Break inner loop\;
        }
        \tcp{slide $j$ if every lane in chunk was less than $s$}
        \If{$L[j+7] < s$}{$j \gets j + 8$; \Continue\;}
        \Break inner loop\;
    }
    \textsc{Scan-Tail}($L, j, s, \mathrm{count}$)
    \Comment*[r]{scalar finish for trailing $< 8$ elements}
}
\Return $\mathrm{count}$\;
\end{algorithm}

The algorithm is the classical ``broadcast-and-compare'' schema
used throughout SIMD intersection literature. Its correctness
follows from the invariant that $j$ is monotonically increased
only after the entire $L[j..j+8]$ window has been ruled out (every
element strictly less than the current $s$); because both $S$ and
$L$ are sorted, no later $s$ can have a match at any index $< j$.

\begin{theorem}[Intersect complexity]
On inputs of sizes $|A|$ and $|B|$ with $|S| = \min(|A|, |B|)$,
$|L| = \max(|A|, |B|)$:
\begin{itemize}
    \item Time: $\Ord{|S| \cdot (|L| / 8 + 1)}$ in the worst case.
        The $/8$ factor is the AVX2 lane width. In practice the
        short-side iteration amortises closer to $\Ord{|S| + |L| / 8}$
        because $j$ slides forward monotonically.
    \item Space: $\Ord{1}$ auxiliary (the register state).
\end{itemize}
\end{theorem}

\begin{theorem}[Runtime safety]
\textsc{Intersect-Sorted-U32-Count} does not invoke undefined
behaviour even when compiled without AVX2, because the unsafe
intrinsic function is only called through \code{is\_x86\_feature\_detected!("avx2")}
at runtime. The compile-time \code{\#[target\_feature(enable =
"avx2")]} gate only enables the intrinsic body inside a function
the scalar path does not call.
\end{theorem}

\begin{proof}[Proof sketch]
The public wrapper is safe. It checks
\code{is\_x86\_feature\_detected!("avx2")} first; only on a hit
does it invoke the unsafe function body. When AVX2 is not reported,
the scalar merge is called and no AVX2 instructions are executed.
The \code{target\_feature} gate means the unsafe function body
itself is compiled with AVX2 codegen regardless, but that's fine
because it's only ever called when the CPU supports it.
\qed
\end{proof}

\section{Roadmap}

Remaining primitives to port from libgraphmatch to
\code{simd\_rust}:
\begin{itemize}
    \item \code{count\_edges\_by\_rel\_type}: scan a
        \code{\&[CompactRelation]} for entries matching a given
        \code{rel\_type\_tag}. Trivial ($\texttt{\_mm256\_cmpeq\_epi8}$
        on the packed tag array).
    \item \code{timestamp\_range\_filter}: mark edges whose
        \code{timestamp} falls in $[t_0, t_1]$. Requires 64-bit
        comparisons --- one extra instruction per lane.
    \item \code{intersect\_sorted\_u32\_out}: the output-materialising
        variant of the current count-only primitive.
\end{itemize}

Once all three land, \textsc{Search-Temporal-Pattern-Simd}'s
libgraphmatch path can be replaced with a pure-Rust dispatch,
retiring the C++ dependency.

\begin{inthecode}
\filepath{graph\_hunter\_core/src/simd\_matcher.rs} --- FFI
bridge.\\
\filepath{graph\_hunter\_core/src/simd\_rust.rs} ---
\textsc{Intersect-Sorted-U32-Count}.\\
\filepath{libgraphmatch/src/simd\_intersect.cpp} --- C++ reference
implementation.\\
\filepath{graph\_hunter\_core/src/graph.rs:689} ---
\textsc{Search-Temporal-Pattern-Simd}.
\end{inthecode}
